
\documentclass{article}
\usepackage{colortbl}
\usepackage{makecell}
\usepackage{multirow}
\usepackage{supertabular}

\begin{document}

\newcounter{utterance}

\centering \large Interaction Transcript for game `ifeval', experiment `ifeval', episode 240 with qwen3.5{-}2b.
\vspace{24pt}

{ \footnotesize  \setcounter{utterance}{1}
\setlength{\tabcolsep}{0pt}
\begin{supertabular}{c@{$\;$}|p{.15\linewidth}@{}p{.15\linewidth}p{.15\linewidth}p{.15\linewidth}p{.15\linewidth}p{.15\linewidth}}
    \# & \multicolumn{2}{c}{Player} && \multicolumn{2}{c}{Game Master} \\
    \hline

    \theutterance \stepcounter{utterance}  
    & & & \multicolumn{4}{p{0.6\linewidth}}{
        \cellcolor[rgb]{0.9,0.9,0.9}{
            \makecell[{{p{\linewidth}}}]{
                \texttt{\tiny{[P1$\langle$GM]}}
                \texttt{Write a casual blog post about how the outer solar system is different from the inner solar system, and what that means for the possibility of life. Wrap your entire response with double quotation marks. Your response should contain 17 or more sentences.} \\
            }
        }
    }
     \\ \\

    \theutterance \stepcounter{utterance}  
    & \multicolumn{4}{p{0.6\linewidth}}{
        \cellcolor[rgb]{0.561,0.737,0.561}{
            \makecell[{{p{\linewidth}}}]{
                \texttt{\tiny{[P1$\rangle$GM]}}
                \texttt{"Imagine standing on a rocky planet far away from the sun, looking out at a sky filled with swirling gas and dust. That is exactly what happens when we look at the outer solar system, which is quite the opposite of the cozy, sun{-}drenched world we are used to. The inner solar system is packed with hot, rocky planets like Mercury, Venus, Earth, and Mars, all orbiting close to the blazing star. In contrast, the outer solar system is dominated by giant, icy gas giants like Jupiter and Saturn, which are much colder and more distant. These distant worlds are covered in thick atmospheres of hydrogen, helium, and methane, creating a very different chemical environment than the rocky surfaces we know. Because the outer planets are so far from the sun, they receive very little heat, allowing them to maintain liquid water on their surfaces for billions of years. This stability of water is crucial, as it suggests that the ingredients necessary for life could have existed in these cold regions long ago. However, the extreme cold and lack of sunlight make it incredibly difficult for complex life forms to evolve or survive in these environments today. The pressure and temperature conditions on the outer planets are so harsh that they are unlikely to support the kind of biological activity we see on Earth. Scientists believe that if life ever arises in the outer solar system, it would likely be microscopic and adapted to survive in freezing conditions. The presence of water ice and organic compounds in the atmospheres of these planets hints at a rich history of chemical reactions that could have led to life. Despite these challenges, the outer solar system remains a fascinating place to study the origins of life and the potential for extraterrestrial biology. We might never find a single living organism there, but the search continues to expand our understanding of the universe. The difference between the inner and outer solar systems highlights how diverse the habitable zones can be across the galaxy. Ultimately, while the inner solar system is our home, the outer solar system offers a unique window into the possibilities of life beyond our planet. It reminds us that life is not limited to the warm, sunlit places we are most familiar with. The journey to explore these distant worlds is both exciting and deeply mysterious. We are still waiting to see if any of these cold giants hold the secrets of life. The answer might just be out there, hidden in the icy dark.} \\
            }
        }
    }
    & & \\ \\

    \theutterance \stepcounter{utterance}  
    & & & \multicolumn{2}{p{0.3\linewidth}}{
        \cellcolor[rgb]{0.9,0.9,0.9}{
            \makecell[{{p{\linewidth}}}]{
                \texttt{\tiny{[GM$|$GM]}}
                {'startend:quotation': {}, 'length_constraints:number_sentences': {'relation': 'at least', 'num_sentences': 17}}
            }
        }
    }
    & & \\ \\

    \theutterance \stepcounter{utterance}  
    & & & \multicolumn{2}{p{0.3\linewidth}}{
        \cellcolor[rgb]{0.9,0.9,0.9}{
            \makecell[{{p{\linewidth}}}]{
                \texttt{\tiny{[GM$|$GM]}}
                \texttt{game\_result = LOSE} \\
            }
        }
    }
    & & \\ \\

\end{supertabular}
}

\end{document}
